\subsection*{Week of 2026-11-09: congruences and number theory}

Second rail: pigeonhole.

\begin{center}\small
\begin{tabular}{@{}c >{\raggedright\arraybackslash}p{3.1cm} >{\raggedright\arraybackslash}p{4.8cm} >{\raggedright\arraybackslash}p{4.8cm}@{}}
\toprule
Rung & Problem & Object and ask & Move \\
\midrule
1 & \putnam{2018 A1} & An equation in positive integers; find all solutions. & Warm: factor and bound. \\
2 & \putnam{2017 B2} & Sums of consecutive integers; a smallest value. & Parity and divisibility of a sum. \\
3 & \seminar{PS 9 \#14} and \seminar{\#15} & Two classic non-divisibility statements. & The order of an element modulo a prime, and the smallest prime factor trick. \\
4 & \putnam{2020 A1} & Integers of a digit pattern; count those divisible by a given number. & Factor the pattern, then orders again. \\
5 & \putnam{2024 A1} & An exponential equation in integers; which exponents. & Reduce modulo small numbers until the cases die. \\
6 & \putnam{2021 A1} & Fixed-length hops in the plane; the minimum number to reach a point. & Tagged pigeonhole and extremal; the number theory of the hop lengths carries it. Cite the rung where you first used a modular obstruction. \\
\midrule
opt & \putnam{2023 B2} & The binary digits of multiples; a minimum. & \\
opt & \putnam{2018 B3} & Integers with linked divisibility conditions. & \\
\bottomrule
\end{tabular}
\end{center}

Seminar warm-ups: \seminar{PS 9 \#1, \#4, \#12, \#31}; \seminar{Pigeonhole \#1, \#6, \#7, \#8, \#9}.
